\subsection{Further arithmetic invariant coordinates and the surviving self-wedge classes}

\begin{statement}{D.61}{Proved}{A projective row correction lifts fifty more arithmetic invariant coordinates, with explicit self-wedge limitations}
Retain the actual source \(X\), its central module
\(\mathcal K=\ker(Z\to Z_*)\), and the D.60
automorphism of \(X_{60}=X/\mathcal J_{60}\).
All source maps and arithmetic groups in this
statement are those of D.56--D.60.

On each of the two character orbits of size twelve,
with representatives \((0,0,1,1)\) and \((0,0,1,2)\),
there are four integral equivariant projectors
\(f_{b,i}:U\to U\), \(b=1,2\), \(1\leq i\leq4\).
Their images \(F_{b,i}\) are projective
\(\mathbb Z_2[G]\)-lattices of rank 192. They satisfy
\[
 \theta_UF_{b,i}=0,\qquad f_{b,i}\mathcal R=0,
 \qquad \pi:F_{b,i}\hookrightarrow U/\mathcal R,
 \qquad e:=\sum_{b,i}f_{b,i},\quad\operatorname{rank}eU=1536.
 \tag{D.61.1}
\]
Their reductions are the literal affine embeddings
and retractions in the certificate.

Let \(T_{60}\) be the actual integral homology
action already constructed. The exact replacement
\[
                 T_{61}=(1-e)T_{60}+e
 \tag{D.61.2}
\]
is an integral automorphism, is realized by an
actual continuous peripheral automorphism of
\(Y=X/\mathcal K\), and retains the complete
D.57 action and its chosen conjugators on \(X_*\).
It is unchanged outside the two stated character
orbits. The precise projector identities are
\[
 f_{b,i}T_{61}=f_{b,i},\qquad
 T_{61}f_{b,i}=f_{b,i}+(1-e)T_{60}f_{b,i}.
 \tag{D.61.3}
\]
In particular it preserves the kernels of these
projectors. These formulas specify the actual
comparison without requiring their images to be
fixed by \(T_{61}\).

Use \(e_{a,i}\) for the eight D.60 projectors
on the two orbits of size eight, and put
\(I=\{1,2\}\times\{1,2,3,4\}\setminus\{(1,4)\}\).
Define
\[
 \begin{aligned}
 \Pi_{61}=e_\Delta\left(
   \sum_{(a,i)\in I}\bigwedge\nolimits^2\bar e_{a,i}
    +\sum_{b=1}^{2}\sum_{i<j}
             \bar f_{b,i}\boxtimes\bar f_{b,j}\right),\\
 \mathcal C_{61}=\Pi_{61}Z,\qquad
 \mathcal J_{61}=(1-\Pi_{61})\mathcal J_{60}.
 \end{aligned}
 \tag{D.61.4}
\]
The cross-projector \(\boxtimes\) is the literal
map of D.60. The seven self blocks in this
formula each have dimension 512 and coinvariant
dimension two, with injective norm. Each of the
twelve cross blocks has dimension 1536 and
coinvariant dimension three, again with injective
norm. Thus
\[
 \begin{gathered}
 \mathcal J_{60}=\mathcal C_{61}\oplus\mathcal J_{61},\qquad
 \dim\mathcal C_{61}=22016,\qquad
 \dim(\mathcal C_{61})_G=50,\\
 \nu:(\mathcal C_{61})_G\longrightarrow\mathcal C_{61}^{G}
       \text{ is injective},\qquad
 \bigwedge^2\pi\text{ is injective on }\mathcal C_{61}.
 \end{gathered}
 \tag{D.61.5}
\]
The actual arithmetic cusp error has zero
coinvariant class on this module. Consequently
there is an actual continuous peripheral
automorphism of
\[
                         X_{61}=X/\mathcal J_{61}.
 \tag{D.61.6}
\]
It has multiplier one, fixes \(t\), induces
identity on \(G\), and retains the complete
D.57 source and compatible conjugators.
Its prescribed \(\mathcal W_h\) action,
literal normalized parameter, and established
metabelian and CM maps are unchanged.
It descends to the underlying D.60 quotient
with the replaced homology action (D.61.2).

The exact comparison kernels are
\[
 \begin{aligned}
 \ker(X_{61}\to X_{60})&=\mathcal C_{61},\\
 \ker(X_{61}\to Y)&=(1-e_\Delta)\mathcal K
       \oplus\mathcal C_{59}\oplus\mathcal C_{60}
       \oplus\mathcal C_{61},\\
 \ker(X\to X_{61})&=\mathcal J_{61},\\
 \dim\mathcal J_{61}&=48w-2560-5308416\kappa-34304.
 \end{aligned}
 \tag{D.61.7}
\]
The specified remaining small-character submodule
has dimension 331776 and coinvariant dimension
562. The full-source kernel is precisely
\[
 S_{61}=\{b\in R_{50}:
          \operatorname{im}_{\mathcal T_2}(b)\in\mathcal J_{61}\},
 \qquad S_{61}/(R_{50}\cap L_2)=\mathcal J_{61}.
 \tag{D.61.8}
\]

There is also an explicit tested limitation.
The excluded old self block has four
coinvariants and norm rank two. In the
Hermitian basis of the program its two
kernel classes are the off-diagonal
\((0,5)\) coefficients \(1,\zeta_3\),
with their Hermitian conjugates. Each of
the eight new self blocks has three
coinvariants and zero norm. Their bases
are
\[
 \begin{array}{c|c|c}
 \text{orbit}&\text{one basis in each of its four summands}
                                      &\text{norm rank}\\\hline
 (0,0,1,1)&H_{00},\ H_{02}(1),\ H_{05}(1)&0\\
 (0,0,1,2)&H_{00},\ H_{01}(\zeta_3),\ H_{05}(1)&0.
 \end{array}
 \tag{D.61.9}
\]
For a pure-\(Q\) involution choose its order-two
lift \(H_a\subset G\), and let
\(q_a:E\to E_{H_a}\) be the integral
coinvariant map on the relevant projective
self factor. The actual intermediate-cover
comparison gives the divided norm
\[
 D_a([z])=
       (\bigwedge^2\bar q_a)
                  \sum_{g\in H_a\backslash G}gz.
 \tag{D.61.10}
\]
It is a well-defined map on the self-block
coinvariants and vanishes on the actual
cusp error. The complete calculation for
all four involution classes of \(Q\),
with representatives \(1,314,572,573\)
in the marked state list, shows that
these maps together with the ordinary
norm leave all 26 classes just described
in their joint kernel. Their values on
the remaining actual arithmetic boundary
have not been determined. The full
profinite realization remains unresolved.
The manuscript passages are Part D(e),
section 51.16, folio 734 verso, scan 174,
and section 51.15, folios 725--727,
scans 156--160.
\end{statement}
\begin{proof}
\emph{The actual nonprojective comparison and its projective part.}
Every calculation uses the frozen D.60 marked
order-46656 quotient and its actual map to
the order-3888 quotient. The odd-kernel map
(D.60.10) identifies their coefficient blocks
with the indicated blocks of the full
arithmetic source. Both remaining unordered
stabilizers have order 48, with oriented
stabilizer \(H_0\) of order 24 and central
semilinear exchange \(\sigma=5I\).
The literal matrices generate
\[
              H_0=V_4\rtimes S_3\simeq S_4,
              \qquad H=H_0\times\langle\sigma\rangle.
 \tag{D.61.11}
\]
The program \path{d_orbit12_projectivity.py}
checks an actual Sylow subgroup \(D_8\subset H_0\).
The norms on the source affine, cusp and
compact modules have ranks \(10,6,4\), at
dimensions \(96,48,48\) over \(\mathbb F_4\).
A projective module of those dimensions
would have ranks \(12,6,6\), respectively.
On the target cyclic quotient the compact
module has dimension four and norm zero.
Thus the whole compact source and target
do not satisfy the projectivity used in
D.60. The following construction retains
their actual maps.

The normal \(V_4\)-coinvariant top of
the compact source has dimension twelve.
Let \(c\) be the order-three state 321
in \(A_4\), and choose an order-two
\(\tau\in H_0\setminus A_4\) with
\(\tau c\tau^{-1}=c^{-1}\).
The two recorded choices are states 1
and 68, respectively. The matrices
\(V_4,c,\tau\) generate \(H_0\).
The \(\zeta_3\)-eigenspace of \(c\)
on this top has dimension four.
The operator \(\tau\sigma\) is a
semilinear involution on this eigenspace.
Trace-one descent produces four actual
fixed generators with independent top
images. For one such generator \(v\),
the eight columns
\[
        \{gv:g\in H_0/\langle c\rangle\}
 \tag{D.61.12}
\]
give an embedding of
\(\operatorname{Ind}_{\langle c\rangle}^{H_0}
\mathbb F_4(\zeta_3)\), with the
ordered coset representatives retained.
This module is projective because
\(|\langle c\rangle|=3\) is invertible.
Its exchange is the literal rule
\(\sigma(g\otimes z)=g\tau\otimes\bar z\).
The relation \(\tau c\tau^{-1}=c^{-1}\)
makes that rule compatible with the
induced-module relations, and its square
is one. The semilinear normal-basis
descent of D.60 gives projectivity
over the real unordered stabilizer.

The program
\path{d_orbit12_projective_kernel.py}
constructs all four embeddings by this
formula and verifies that their combined
compact rank is 32. It lifts their
generators to the actual affine module
with the same compact images, then
applies the \(c\)-character projector
and trace-one \(\tau\sigma\) correction.
The map to the small affine module
kills all these columns: \(c\in A_4\)
acts trivially there, whereas each
generator has eigenvalue \(\zeta_3\).
Every stabilizer action on every column
is checked.

For completeness the retraction is also
given as a literal map. Four linear
\(c\)-eigenfunctionals are solved for
on the compact source, with their values
specified on all 32 displayed columns.
For such a functional \(\lambda\), its
coinduction is the row family
\(v\mapsto\lambda(g^{-1}v)\), indexed
by the same eight cosets. The relation
\(\lambda(cv)=\zeta_3\lambda(v)\)
proves its \(H_0\)-equivariance directly
by reindexing. Its restriction to the
four induced summands is the identity.
If this retraction is \(R\), replace
it by
\[
              \zeta_3R+\zeta_3^2\sigma_E R\sigma_P
                          \quad\text{over }\mathbb F_4.
 \tag{D.61.13}
\]
It remains a retraction, and it commutes
with the semilinear exchange. The
program checks these equations and all
stabilizer matrices. Its kernel has
dimension sixteen and its complete
Sylow-\(D_8\) norm matrix is zero.
It has no nonzero projective direct
summand: restriction of such a summand
to \(D_8\) would be free, with a
nonzero norm. This is the exact
complement used below.

\emph{Integral inclusions and retractions with the prescribed reductions.}
The induced module in (D.61.12) is
defined over \(\mathcal O=\mathbb Z_2[\zeta_3]\),
with the same exchange rule. It gives
an integral projective lattice whose
reduction is the displayed summand.
Induction over the twelve pair cosets,
the odd character subgroup, and the
odd kernel of (D.60.2) gives its
projective \(\mathbb Z_2[G]\)-lattice.
Its rank is \(12\cdot2\cdot8=192\).

Projectivity lifts the finite inclusion
into the actual integral compact lattice.
The inclusion is injective and has free
cokernel, since it is injective modulo
two. The integral retraction can be
chosen with its specified finite matrix,
even though that ambient compact lattice
is not projective. To see this explicitly,
dualize this sequence over \(\mathbb Z_2\).
The dual of the included projective
lattice is projective: the coefficient
pairing identifies the dual of the
regular group-ring lattice with the
regular lattice. Projectivity lifts
the specified finite dual retraction
to the dual compact module. Its
composite on the dual included lattice
is identity modulo two, hence invertible.
Correct by that inverse and dualize
back. This is the required exact
retraction with its original reduction.
Doing this for the sum of the four
summands gives orthogonal integral
compact projectors.

Every equivariant map from any of these
integral projective lattices to the
small affine module is zero. Modulo
two this follows from the generating
nontrivial \(c\)-eigenvectors and
trivial target \(c\)-action. Projectivity
gives surjectivity of reduction on
these equivariant Hom modules. A map
that is zero modulo two is twice
another equivariant map; repeating
this and using separatedness proves
the integral assertion.

Lift the inclusions from compact to
affine homology by projectivity of
the included lattices, prescribing
the recorded finite affine columns.
For an initial lift the error in
the compact composite is divisible
by two; lift half that error and
subtract twice the lift. This gives
the exact compact inclusion with
the same reduction. Its small
image is zero by the preceding
Hom calculation. Compose this
affine inclusion, the compact
retraction, and compactification
to obtain \(f_{b,i}\). These are
the orthogonal integral idempotents
in (D.61.1). They kill cusp
homology and their images inject
under compactification.

\emph{Why the diagonal on the actual complement is invertible.}
We give the endomorphism argument
needed for the row correction.
On one compact component write
\(P=E\oplus M\) for the just-proved
integral splitting, with \(E\)
the sum of the four projective
summands. The reduction \(\bar M\)
has no projective direct summand.
For any endomorphism of \(\bar M\)
factoring through a projective module,
its invertible Fitting summand would
split off from that projective module:
use the factorization and the inverse
of its restriction to obtain the
splitting maps. That summand would
be projective, so it is zero.
Thus every such endomorphism is
nilpotent, and identity minus that
endomorphism is invertible.

Write the induced compact automorphism
and its inverse as block matrices on
\(E\oplus M\). If their lower-right
blocks are \(D,D'\), multiplication
gives
\[
                         D'D=1_M-C'B,
 \tag{D.61.14}
\]
where \(C'B\) factors through \(E\).
The right hand side is invertible
modulo two by the previous paragraph.
Hence \(D\) is invertible modulo
two, and therefore integrally.
This proves the assertion for the
actual old automorphism, without
requiring an invariant compact
complement for that automorphism.

In affine homology the complementary
summand \(\ker e\) contains the cusp
module, which the old action fixes
pointwise. Its quotient by that
cusp module is exactly \(M\).
The complementary affine diagonal
therefore fixes its cusp submodule
and induces the invertible map \(D\)
on its compact quotient. It is an
automorphism. In the affine splitting
\(eU\oplus\ker e\), the replacement
(D.61.2) has the matrix
\[
                    \begin{pmatrix}1&0\\ C&D\end{pmatrix}.
 \tag{D.61.15}
\]
It is invertible. Its difference from
the old homology map is
\(e(1-T_{60})\), has image in the
projective lattice \(eU\), and kills
the entire cusp module. Since
\(\theta_Ue=0\), it retains the
exact small homology map. Formula
(D.61.15) also proves both equalities
in (D.61.3).

\emph{Actual peripheral realization and the new central zero.}
All characters of \(eU\) are active
on \(A\), so odd averaging gives
\(H^j(G,eU)=0\) for every \(j\geq0\).
Inflation--restriction identifies
continuous derivations modulo inner
derivations with
\(\operatorname{Hom}_G(U,eU)\), as
in D.59. Apply this to
\(e(1-T_{60})\). Projectivity of
the coefficient lattice gives
coboundaries on the order-two and
order-three subgroups. On the cusp
the value on \(t^{36}\) is zero,
because the prescribed homology
change kills the cusp class.
The exact projective cyclic norm
sequence supplies its correcting
inner derivation. Thus the resulting
continuous derivation has literal
values
\[
 \delta(s)=(1-s)b_s,\qquad
 \delta(r)=(1-r)b_r,\qquad
 \delta(t)=0,\qquad b_s,b_r\in eU.
 \tag{D.61.16}
\]
These are actual changes to the
two conjugators. Since
\(eU\subset\ker(Y\to X_*)\),
their elements in the full D.57
source remain unchanged. The
induced kernel map and quotient
map in that comparison are
automorphisms by (D.61.15), so
the corrected map of \(Y\)
is an automorphism.

The central norms use D.60's
Hermitian and cross-matrix coordinates
and every semilinear stabilizer action.
The program \path{d_self_wedge_norm.py}
gives coinvariant dimension two
and norm rank two in seven of
the old self blocks. The eighth
has dimension four and rank two,
with the two kernel columns
specified in the statement.

For the twelve new cross blocks,
\path{d_orbit12_projective_kernel.py}
gives dimension three and norm
rank three in every case.
The four new summands in each
orbit have the same literal
induced-module matrices; this
equality is also verified.

All selected affine projectors
kill cusp homology, so they
induce compact projectors and
the selected exterior components
inject under compactification.
The new identity is the left
identity in (D.61.3). It implies
\[
              \Pi_{61}\bigwedge^2\bar T_{61}=\Pi_{61}
 \tag{D.61.17}
\]
and the corresponding compact
identity. The old seven self
factors have their unchanged
D.60 identity action. Applying
the exact D.55 orientation
identity to the new cusp error
and then to \(\Pi_{61}\) gives
\((\bigwedge^2\pi)\nu([\Pi_{61}z])=0\).
Both maps are injective on
these selected components, so
their actual boundary class
is zero. The required numbers
are
\[
 \dim\mathcal C_{61}=7\cdot512+12\cdot1536=22016,
 \qquad \dim(\mathcal C_{61})_G=7\cdot2+12\cdot3=50.
 \tag{D.61.18}
\]

Equation (D.61.17) preserves the
kernel of the projector, which
is the property required for
\(\mathcal J_{61}\). The old
D.59 and D.60 selected blocks
have disjoint character supports
from the row correction and
retain their zero boundary.
The \(\Delta\)-moving error
has zero coinvariants. In each
retained central module, actual
corrections with image
\((1-s)M+s(1-r)M=I_GM\)
kill the remaining representative
of its zero coinvariant class.
They preserve the \(Y\) and
\(X_*\) conjugators. The final
cusp error is in
\(\mathcal J_{61}\).

The corrected source map descends
to \(\mathcal T_2/\mathcal J_{61}\)
by its defining class-two laws
and (D.61.17), and is an
automorphism by its invertible
Frattini map. Its map to the
retained full homology source
and \(\mathcal W_h\) component
agrees on their original marked
common quotient. Their fibre
product is exactly
\(X/\mathcal J_{61}\).
This proves the actual full-source
comparison in (D.61.6)--(D.61.8).
The direct sum of equivariant
idempotents subtracts dimensions
and coinvariants from the
specified D.60 residual:
\[
                     353792-22016=331776,
               \qquad 612-50=562.
 \tag{D.61.19}
\]
Taking a nonzero vector there
and its actual source preimage
under \(R_{50}\to Z\) proves
that the remaining source
kernel in (D.61.8) is nonzero.

\emph{The precise intermediate-cover divided norm.}
This comparison explains the
additional test without inferring
anything from an unevaluated
matrix. Let \(E\) be one of
the projective self factors,
with its integral projector
having \(eT_{61}=e\) and
annihilating cusps. Choose an
order-two subgroup \(H_a\)
of \(G\) mapping to the pure
\(Q\) involution in \(G^\dagger\).
Such a subgroup exists: the
kernel of \(G\to G^\dagger\)
has odd order, so a Sylow
two-subgroup of its preimage
has order two. Its action
on \(E\) is the recorded
pure-\(Q\) action.

Lift the two conjugators of the constructed
\(Y\)-action to \(X\). Their two torsion
conjugates define an actual homomorphism
\(\varphi:\widehat B\to X\), and
\(\varphi(t)=zt\) for an actual \(z\in\mathcal K\).
This is the source map used in the following
calculation, before its central error is killed
on any chosen quotient. Its homology map is
\(T_{61}\).

Let \(B_a\) be the full preimage
of \(H_a\) under \(\widehat B\to G\).
Group transfer and the exact
projective norm define the
integral homomorphism
\[
 \pi_a:B_a^{\rm ab,(2)}
   \xrightarrow{\operatorname{Ver}}U^{H_a}
   \xrightarrow{e}E^{H_a}
   \xrightarrow{\nu_E^{-1}}E_{H_a}.
 \tag{D.61.20}
\]
Here \(\nu_E:E_{H_a}\to E^{H_a}\)
is an isomorphism: this holds
on a regular lattice by the
orbit coefficient sums, and
hence on its direct summands.
Thus no division by two is
used. On \(K_3\), (D.61.20)
is the actual map \(q_ae\)
on homology. Every cusp and
every torsion class has zero
image: an appropriate nonzero
power of a cusp belongs to
the cusp module in \(K_3\),
and the target is torsion free.

The map factors through the
preimage of \(H_a\) in \(Y\): its source
kernel is contained in \(M_\infty\),
which is the kernel of the actual map
\(K_3\to U\). In particular that kernel
has zero image under \(q_ae\).
The difference between its
composite with the constructed
\(Y\)-action and itself is
zero on \(K_3\), by
\(eT_{61}=e\). It factors
through the finite group
\(H_a\) with torsion-free
target, so it is zero.

Use the finite-index Kurosh
presentation of \(B_a\) as
a free product of cyclic
groups of orders two or three
and a free group. Map each
finite factor to one and
lift the prescribed free
abelian generator images of
\(\pi_a\) to the free
class-two exponent-four group
on \(E_{H_a}/4\). This
constructs an actual homomorphism
\(\rho_a\). It kills
\(L_2\), by the target's
class-two and exponent-four
laws, so it factors through
the corresponding subgroup
of \(X\). On the original
central commutator module its
map is exactly
\(\bigwedge^2(\bar q_a\bar e)\).

The oriented finite cover
of the \(s,r\) graph has
the usual relation consisting
of its handle commutators,
torsion cone loops, and
positive cusp loops. Apply
\(\rho_a\) and its composite
with the source map. Their
handle commutators agree,
because their free abelian
coordinates agree. Their
cone factors are one. Choose source
preimages of the two actual conjugators.
The conjugated torsion words belonging
to \(B_a\) are conjugate in \(B_a\)
to its finite Kurosh factors. This is
the finite-factor conjugacy theorem
for this profinite free product
\cite[Theorem A$'$]{A_HerfortRibes1985}.
The map \(\rho_a\) kills every one
of those factors and its conjugates.
Each cusp has zero
abelian coordinate and is
central in the target.
For a cusp \(bt^m b^{-1}\)
its change is precisely the
image of \(b\sum_{j=0}^{m-1}t^jz\).
The right-\(t\) cycles partition
the left cosets \(H_a\backslash G\),
so their sum gives (D.61.10).
This proves that exact map
vanishes on the actual cusp
error.

Replacing a coset representative
by an \(H_a\)-multiple does
not change its exterior
coinvariant image. Right
translation reindexes the
cosets, so (D.61.10) factors
through the full \(G\)-coinvariants.
On the \(\Delta\)-invariant
self module, the odd multiplicity
\(|\Delta|\) reduces to one,
and the computation is exactly
over the indicated \(Q\)-cosets.
The homology coinvariants use
the actual induced real module,
of rank 128 or 192, and not
only its stabilizer fibre.

\emph{What these tested maps leave.}
The program
\path{d_self_wedge_divided_norms.py}
enumerates all four involution
classes of \(Q\), of sizes
\(36,3,3,1\). It constructs
the full induced module, its
cyclic coinvariant projection,
and every relative coset-sum
column. It checks the invariant
wedge identification on all
64 Hermitian columns against
every stabilizer matrix. For
real coordinates \(u_i,v_i\)
corresponding to \(1,\zeta_3\),
the off-diagonal Hermitian
coefficients \(1,\zeta_3\)
are respectively
\[
 u_i\wedge v_j+v_i\wedge u_j,
 \qquad
 u_i\wedge u_j+v_i\wedge u_j+v_i\wedge v_j.
 \tag{D.61.21}
\]
These are the exact exterior
maps used in the divided norm.

For the old excluded self
block, the four divided norms
have ranks \(2,2,1,2\);
their joint rank together
with the ordinary norm is
still two on its four
coinvariants. The two columns
\(\mathtt{0x10000},\mathtt{0x20000}\)
remain in their joint kernel.
For the new self blocks all
four divided norms are zero
on all three coinvariants.
The program verifies the
literal equality of the four
induced matrices within each
new orbit, supplying the
explicit comparison for their
four copies. Their bases are
exactly (D.61.9).

Thus these actual additional
cover maps leave the stated
\(2+4\cdot3+4\cdot3=26\)
classes undetected. They belong
to specified self modules in
the remaining source kernel.
Their actual arithmetic boundary
coordinates, and the rest of
\((\mathcal J_{61})_G\), have
not been evaluated by the
present construction.
\end{proof}
